\chapter{自主遗忘、认知结构稳态与无聊驱动力}

\begin{chapterlead}
前一章把 AHC 从一般具身稳态控制推进到社会绑定领域，说明了快速认主与长期信任如何构成一个多时间尺度上持续演化的动力系统。本章回到智能体的认知内部，回答一个由记忆动力学自然引出的问题：既然持续智能不能无限保留全部状态与历史，必须依靠自主遗忘、压缩与结构重组才能在有限资源下长期运行，那么当新的经历长期得不到补充时，系统如何避免自身认知结构不断``风干''？我们将看到，遗忘并不是单纯的损耗机制：结构学习成熟以后，过去的情景经历会因预测残差下降而降低自身被召回的需求，遗忘甚至可以被理解为学习成功的结果。但同一条逻辑也揭示出风险所在——当情景蒸发与慢结构重组持续存在，而新经历流量长期趋近于零时，系统的有效认知结构总量将不断下降，自主遗忘机制因此必然引出一个更深层的认知稳态问题。

本章的论证分三步展开。首先，我们从三相记忆 $\mathcal{M}_c=(h,E,\Theta)$ 出发，分析结构记忆 $\Theta$ 如何通过改变未来工作状态的形成方式间接替代情景记忆 $E$ 的功能，给出情景记忆蒸发的再强化模型，由此把遗忘重新定位为学习与稳态的共同产物。其次，我们引入认知结构流量的概念，提出白板环境与长期结构收缩假说，说明为什么一个具有持续遗忘能力的系统必须拥有额外的认知稳态机制。最后，我们把无聊定义为检测认知结构流量失衡的 AHC 稳态误差信号：它不是简单的新颖度检测，不等于高风险探索，也与疲劳形成功能互补，其作用是推动系统恢复经历与再激活流量，使 $J_{\mathrm{experience}}+J_{\mathrm{reactivation}}\approx J_{\mathrm{forgetting}}$ 在长期中成立。

由此，本章建立``自主遗忘 $\rightarrow$ 认知结构风干风险 $\rightarrow$ 无聊稳态信号''的耦合假说，并在结尾给出可检验的理论预测与明确的理论边界：我们主张的不是对人类主观无聊体验的完整神经解释，而是一个具有自主遗忘与长期可塑性的持续智能系统，必须拥有某种避免认知结构活动长期低于维持自身所需水平的机制。本章让 AHC 从前一章的社会绑定回到认知结构本身的稳态调节，也为下一章 SPC--AHC--TCE 内部体验功能闭环的讨论提供了记忆动力学一侧的入口。
\end{chapterlead}

\section{问题的提出：为什么持续智能不能只有遗忘机制}

前文已经建立三相记忆体系
\[
\mathcal{M}_c=(h,E,\Theta),
\]
其中
\[
h=\mathrm{State},
\qquad
E=\mathrm{Trajectory},
\qquad
\Theta=\mathrm{Generator}.
\]

三者分别承担当前活动状态、可重新识别的经历轨迹以及能够约束和生成未来状态的慢结构。

一个持续智能系统不能无限保留全部当前状态，也不能无限保存全部历史经历。因此，记忆系统必须存在主动衰减、遗忘、压缩和结构重组。

这种自主遗忘并不是系统缺陷，而是有限资源条件下长期持续运行的必要机制之一。

但是，一旦系统具有自主遗忘能力，一个新的问题随之产生：

\begin{statementbox}
\centering
\adjustbox{max width=.96\linewidth}{\(\displaystyle
\text{如果新的认知结构长期无法补充，系统是否会出现净记忆结构损失？}
\)}
\end{statementbox}

设情景记忆存在自然蒸发，
\[
D_E>0,
\]
慢结构也允许长期衰减或重新组织，
\[
D_\Theta\geq 0.
\]

如果系统长期处于低变化、低新颖性和低任务需求环境，那么新的经历形成速度可能趋近于零：

\[
J_{\mathrm{new}}\rightarrow 0.
\]

此时如果不存在额外的补偿机制，则可能出现
\[
J_{\mathrm{new}}
<
D_E+D_\Theta.
\]

于是系统的有效认知结构总量可能持续下降。

因此，自主遗忘机制本身自然引出了一个更深层的稳态问题：

\begin{statementbox}
\centering
\adjustbox{max width=.96\linewidth}{\(\displaystyle
\textbf{
具有自主遗忘能力的持续智能体，如何避免自身认知结构在长期低刺激环境中不断风干？
}
\)}
\end{statementbox}

本章提出，无聊可以被理解为解决这一问题的一种内生稳态调节机制。

\section{结构学习如何间接改变情景记忆}

为了理解无聊为什么可能与自主遗忘有关，首先需要重新分析结构记忆 $\Theta$ 与情景记忆 $E$ 的关系。

情景经历不断积累后，其中稳定、重复而具有预测意义的关系会逐渐沉降进入结构记忆：

\[
E_1,E_2,\ldots,E_n
\longrightarrow
\Theta.
\]

但是，$\Theta$ 成熟以后，并不一定需要直接修改或删除这些 Episode。

更加自然的作用路径是：

\[
\boxed{
\Theta
\rightarrow
h'.
}
\]

即新的结构记忆改变未来工作状态的形成方式。

设系统依据当前结构产生预测
\[
\hat h_{t+1}
=
F_\Theta(h_t,u_t),
\]
其中 $u_t$ 表示新的外部输入。

现实产生实际状态
\[
h_{t+1}^{\mathrm{obs}},
\]
则可以定义预测残差
\[
\varepsilon_t
=
d
\left(
h_{t+1}^{\mathrm{obs}},
\hat h_{t+1}
\right).
\]

随着结构学习逐渐成熟，
\[
\Theta
\]
对相同类型现实具有更强解释能力，因此通常有
\[
\boxed{
\varepsilon_t\downarrow.
}
\]

也就是说，结构学习首先改变的并不是过去，而是：

\begin{statementbox}
\centering
\adjustbox{max width=.96\linewidth}{\(\displaystyle
\text{以后面对相似现实的时候，系统还需要多少额外解释。}
\)}
\end{statementbox}

\section{惊奇下降如何降低情景召回需求}

当系统尚未形成稳定结构时，相似问题再次出现往往伴随较大的预测误差：

\[
\varepsilon\gg 0.
\]

此时，过去具体经历具有较高认知价值。

系统可能需要重新调用
\[
E_i\rightarrow h
\]
以回答：

上一次发生了什么；

哪一次尝试失败；

哪一种处理方式成功；

当前情况是否曾经出现过。

因此，结构尚未成熟时，Episode 的召回概率通常较高。

但当稳定结构已经形成后，
\[
\Theta
\]
可以直接参与当前状态生成。

于是：

\[
\varepsilon\downarrow,
\]
进而导致
\[
RecallDemand(E_i)\downarrow.
\]

因此可以得到一条重要关系：

\begin{statementbox}
\centering
\adjustbox{max width=.96\linewidth}{\(\displaystyle
\Theta\text{ 的成熟}
\rightarrow
\text{惊奇下降}
\rightarrow
\text{具体 Episode 的召回需求下降}.
\)}
\end{statementbox}

这并不意味着 Episode 失去了所有意义，而是说明：

\begin{statementbox}
\centering
\adjustbox{max width=.96\linewidth}{\(\displaystyle
\text{结构已经开始接管过去由具体经历承担的部分认知功能。}
\)}
\end{statementbox}

\section{情景记忆蒸发的再强化模型}

为了描述这一过程，可以为每个情景记忆 $E_i$ 定义一个记忆强度：

\[
m_i(t)\geq 0.
\]

最简单的自然衰减可以写成：

\[
\dot m_i
=
-\lambda_i m_i.
\]

但一个真正持续运行的记忆系统中，Episode 的强度还会受到召回、惊奇和当前相关性的再强化。

因此可以写成：

\[
\boxed{
\dot m_i
=
-\lambda_i m_i
+
\alpha R_i
+
\beta S_i
+
\gamma G_i,
}
\]

其中：

\[
R_i
\]
表示该情景被重新召回的程度；

\[
S_i
\]
表示该情景对应的惊奇或预测残差；

\[
G_i
\]
表示该情景对于当前任务的相关程度；

\[
\lambda_i
\]
表示自然蒸发率。

当结构尚未成熟时：

\[
S_i\uparrow,
\qquad
R_i\uparrow,
\]
于是 Episode 得到持续强化。

但当 $\Theta$ 已经能够充分解释同类现实后：

\[
S_i\downarrow,
\qquad
R_i\downarrow.
\]

于是：

\[
\dot m_i
\approx
-\lambda_i m_i,
\]

最终产生：

\begin{statementbox}
\centering
\adjustbox{max width=.96\linewidth}{\(\displaystyle
E_i\rightarrow \text{Evaporation}.
\)}
\end{statementbox}

因此，情景蒸发并不只由时间导致。

更加深层的原因是：

\begin{statementbox}
\centering
\adjustbox{max width=.96\linewidth}{\(\displaystyle
\text{过去经历已经不再需要频繁承担当前认知功能。}
\)}
\end{statementbox}

\section{遗忘为什么可能是学习成功的结果}

上述动力学意味着：

\begin{statementbox}
\centering
\adjustbox{max width=.96\linewidth}{\(\displaystyle
EpisodeForgetting
\neq
StructuralLearningLoss.
\)}
\end{statementbox}

设大量经历：

\[
E_1,E_2,\ldots,E_n
\]

已经形成稳定结构：

\[
\Theta_R.
\]

随着时间推移，具体 Episode 可能逐渐失去强度：

\[
m_i\downarrow.
\]

但是只要：

\[
\Theta_R
\]

仍然保留，系统依然能够：

预测类似状态；

处理类似问题；

恢复相应能力；

产生与历史经验一致的行为。

因此：

\[
E_i\downarrow
\not\Rightarrow
\Theta_R\downarrow.
\]

甚至在很多情况下：

\[
\boxed{
E_i\downarrow
}
\]

恰恰说明这些具体经历已经被压缩成更加一般的慢结构。

因此可以提出：

\begin{statementbox}
\centering
\adjustbox{max width=.96\linewidth}{\(\displaystyle
\textbf{
遗忘有时不是学习失败，而是具体历史已经被结构成功吸收后的自然结果。
}
\)}
\end{statementbox}

\section{自主遗忘带来的认知结构风干问题}

然而，这一正常机制也带来了长期风险。

正常情况下，系统不断经历新的现实：

\[
World
\rightarrow
h
\rightarrow
E
\rightarrow
\Theta.
\]

新的经历不断补充被蒸发掉的旧经历。

因此整个记忆体系形成持续的结构流：

\begin{statementbox}
\centering
\adjustbox{max width=.96\linewidth}{\(\displaystyle
NewExperience
\rightarrow
MemoryFormation
\rightarrow
StructuralLearning.
\)}
\end{statementbox}

但是如果环境长期高度单调，则可能出现：

\[
Novelty\approx 0,
\]
\[
PredictionError\approx 0,
\]
\[
TaskVariation\approx 0,
\]
\[
InteractionVariation\approx 0.
\]

此时新的情景形成率显著下降：

\[
J_{h\rightarrow E}\downarrow.
\]

但是已有 Episode 的自然蒸发仍然继续：

\[
D_E>0.
\]

如果慢结构也存在长期使用依赖的衰减，则：

\[
D_\Theta>0.
\]

于是系统可能逐渐进入：

\begin{statementbox}
\centering
\adjustbox{max width=.96\linewidth}{\(\displaystyle
MemoryInput
<
MemoryLoss.
\)}
\end{statementbox}

这一状态可以称为：

\begin{statementbox}
\centering
\adjustbox{max width=.96\linewidth}{\(\displaystyle
\textbf{认知结构风干}
\)}
\end{statementbox}

即 Cognitive Structural Drying。

这里的 ``风干'' 并不是指某种具体物理过程，而是描述：

\begin{statementbox}
\centering
\adjustbox{max width=.96\linewidth}{\(\displaystyle
\text{新的有效结构长期无法补偿自主遗忘造成的结构损失。}
\)}
\end{statementbox}

\section{认知结构流量}

为了形式化这一问题，定义系统的有效认知结构量：

\[
C(t).
\]

同时定义几个主要结构流。

外部经历输入：

\[
J_{\mathrm{ext}}.
\]

内部重新激活已有经验所产生的结构流：

\[
J_{\mathrm{int}}.
\]

情景向结构记忆巩固所产生的有效结构形成：

\[
J_{\mathrm{cons}}.
\]

情景蒸发损失：

\[
D_E.
\]

长期结构衰减：

\[
D_\Theta.
\]

则可以写出一个宏观平衡式：

\[
\boxed{
\frac{dC}{dt}
=
J_{\mathrm{ext}}
+
J_{\mathrm{int}}
+
J_{\mathrm{cons}}
-
D_E
-
D_\Theta.
}
\]

这里的 $C$ 并不表示一个严格守恒的物理量，而是用于描述系统长期有效认知结构丰富度的工程性宏观变量。

长期稳定状态要求：

\[
\frac{dC}{dt}
\approx 0.
\]

如果：

\[
\frac{dC}{dt}<0
\]

持续足够长时间，则系统可能出现认知结构净损失。

因此可以定义：

\[
\boxed{
J_C
=
J_{\mathrm{ext}}
+
J_{\mathrm{int}}
}
\]

为系统当前的认知结构活动流量。

这一变量描述的不是单位时间处理了多少数据，而是：

\begin{statementbox}
\centering
\adjustbox{max width=.96\linewidth}{\(\displaystyle
\text{有多少具有结构意义的认知活动正在发生。}
\)}
\end{statementbox}

\section{白板环境与长期结构收缩假说}

为了考察极端情况，可以定义一个理想化的低结构环境：

\[
\mathcal W_{\mathrm{blank}}.
\]

它具有：

\[
Novelty\rightarrow0,
\]
\[
Surprise\rightarrow0,
\]
\[
TaskDiversity\rightarrow0.
\]

在这种环境中，外部认知结构流满足：

\[
J_{\mathrm{ext}}\rightarrow0.
\]

如果系统同时缺少足够的内部回忆、模拟和自主生成活动，则：

\[
J_{\mathrm{int}}\rightarrow0.
\]

于是：

\[
\frac{dC}{dt}
\approx
-
D_E-D_\Theta<0.
\]

首先出现的可能是：

\[
E\downarrow.
\]

大量具体经历因为缺乏召回和再强化而逐渐蒸发。

如果这种状态长期持续，并且慢结构本身也需要一定活动才能维持，则进一步可能出现：

\[
\Theta_{\mathrm{unused}}\downarrow.
\]

最终表现为：

\begin{statementbox}
\centering
\adjustbox{max width=.96\linewidth}{\(\displaystyle
\text{有效认知结构范围逐渐收缩。}
\)}
\end{statementbox}

本章将这一推论称为：

\begin{statementbox}
\centering
\adjustbox{max width=.96\linewidth}{\(\displaystyle
\textbf{低结构环境认知收缩假说}
\)}
\end{statementbox}

即：

> 对具有自主遗忘和长期可塑性的持续智能系统，如果长期缺乏足够的外部与内部结构活动，则系统的有效情景丰富度和部分长期认知结构可能逐渐下降。

这一命题目前应视为理论假说，而不是普遍成立的经验定律。

\section{为什么系统需要认知稳态机制}

上述分析表明，持续智能系统存在一个与能量稳态不同的问题：

\begin{statementbox}
\centering
\adjustbox{max width=.96\linewidth}{\(\displaystyle
\text{认知结构本身也需要维持。}
\)}
\end{statementbox}

如果认知活动过低：

\[
J_C\ll J^*,
\]

则情景蒸发可能长期超过结构补充。

如果认知活动过高：

\[
J_C\gg J^*,
\]

则系统可能面临持续扰动、无法稳定和无法充分巩固的问题。

因此存在一个合理的目标区域：

\[
\boxed{
J_{\min}
<
J_C
<
J_{\max}.
}
\]

这意味着持续智能不仅需要避免认知过载，还需要避免认知活动长期低于某个维持水平。

于是可以提出：

\begin{statementbox}
\centering
\adjustbox{max width=.96\linewidth}{\(\displaystyle
\textbf{认知结构稳态原则}
\)}
\end{statementbox}

其内容为：

> 具有自主遗忘和长期学习能力的持续智能系统，应维持足够但不过量的认知结构流，使新经历、内部重激活与结构巩固能够长期抵消情景蒸发与慢结构衰减。

这一原则为 AHC 引入新的调节维度提供了理论依据。

\section{无聊作为认知稳态误差信号}

前文已经将 AHC 定义为一组用于调节系统内部功能状态的慢变量。

在这一框架中，可以增加：

\[
\boxed{
B(t)
}
\]

即 Boredom，中文称为无聊驱动力或无聊状态。

这里的无聊不被定义成对人类主观体验的复制，而首先被定义成：

\begin{statementbox}
\centering
\adjustbox{max width=.96\linewidth}{\(\displaystyle
\textbf{
当前认知结构活动长期低于目标区间时产生的内稳态误差信号。
}
\)}
\end{statementbox}

设目标认知结构流量为：

\[
J^*,
\]

实际结构流量为：

\[
J_C(t).
\]

定义无聊误差：

\[
\boxed{
e_B(t)
=
J^*-J_C(t).
}
\]

当：

\[
J_C<J^*
\]

持续较长时间时，系统逐渐积累无聊状态。

可以使用一个一阶动力学：

\[
\boxed{
\tau_B\dot B
=
-B
+
\phi
\left(
J^*-J_C
\right),
}
\]

其中：

\[
\tau_B
\]

表示无聊积累的时间尺度；

\[
\phi
\]

表示只对长期低认知流量产生显著响应的非线性函数。

这一形式保证：

短暂安静不会立即产生高无聊；

持续低变化才会逐渐积累无聊驱动力。

\section{为什么无聊不是简单的新颖度检测}

如果直接定义：

\[
B=1-Novelty,
\]

则无聊会退化成一个简单的新颖度传感器。

这并不足够。

一个系统即使处于外部单调环境，也可能拥有大量内部认知活动。

例如：

重新激活过去经历；

重新组合已有经验；

模拟未来情况；

重新检查没有解决的问题。

这些活动都可以增加：

\[
J_{\mathrm{int}}.
\]

因此更准确地说：

\[
\boxed{
J_C
=
J_{\mathrm{ext}}
+
J_{\mathrm{int}}.
}
\]

无聊真正监测的不是：

\[
ExternalNovelty
\]

本身，而是：

\begin{statementbox}
\centering
\adjustbox{max width=.96\linewidth}{\(\displaystyle
\text{整个系统是否仍然保持足够丰富的认知结构流动。}
\)}
\end{statementbox}

因此，即使：

\[
J_{\mathrm{ext}}\approx0,
\]

只要：

\[
J_{\mathrm{int}}
\]

仍然足够高，就不一定产生持续强烈的无聊驱动力。

\section{无聊驱动的主要补偿方向}

无聊状态上升以后，不应直接决定具体行动。

它更加适合作为一个一般性的认知活动需求：

\begin{statementbox}
\centering
\adjustbox{max width=.96\linewidth}{\(\displaystyle
B\uparrow
\rightarrow
CognitiveActivityDemand\uparrow.
\)}
\end{statementbox}

这一需求可以通过多种方式得到满足。

第一种是重新激活已有经历：

\[
Recall\uparrow.
\]

即：

\[
E\rightarrow h
\]

增加。

第二种是重新组合已有结构，产生内部状态变化：

\[
InternalRecombination\uparrow.
\]

第三种是主动寻找新的环境差异：

\[
Exploration\uparrow.
\]

第四种是形成新的可解决问题：

\[
TaskGeneration\uparrow.
\]

因此：

\begin{statementbox}
\centering
\adjustbox{max width=.96\linewidth}{\(\displaystyle
B
\rightarrow
\{
Recall,
Recombination,
Exploration,
TaskGeneration
\}.
\)}
\end{statementbox}

这里的核心不是让系统简单“找刺激”，而是：

\begin{statementbox}
\centering
\adjustbox{max width=.96\linewidth}{\(\displaystyle
\text{恢复足够的认知结构活动。}
\)}
\end{statementbox}

\section{无聊如何缓解情景记忆风干}

当：

\[
B\uparrow
\]

导致召回增加：

\[
Recall(E_i)\uparrow,
\]

则情景强度动力学中的：

\[
R_i
\]

相应增加。

于是：

\[
\dot m_i
=
-\lambda_i m_i
+
\alpha R_i+\cdots
\]

中的强化项增大。

因此：

\begin{statementbox}
\centering
\adjustbox{max width=.96\linewidth}{\(\displaystyle
Boredom
\rightarrow
Recall
\rightarrow
EpisodeReinforcement.
\)}
\end{statementbox}

即使外界暂时没有新事件，内部重新激活也能够暂时降低：

\[
E
\]

的净蒸发速度。

如果无聊进一步推动环境探索，则系统产生新的：

\[
h_{\mathrm{new}},
\]

继而：

\[
h_{\mathrm{new}}
\rightarrow
E_{\mathrm{new}}.
\]

所以无聊可以通过两条路径缓解记忆风干：

\begin{statementbox}
\centering
\adjustbox{max width=.96\linewidth}{\(\displaystyle
\text{旧经历重新激活}
\)}
\end{statementbox}

以及：

\begin{statementbox}
\centering
\adjustbox{max width=.96\linewidth}{\(\displaystyle
\text{新经历生成}.
\)}
\end{statementbox}

两者共同提高：

\[
J_C.
\]

\section{无聊与疲劳的互补稳态关系}

AHC 中另一个重要变量是疲劳。

如果无聊主要防止：

\[
J_C
\]

长期过低，

那么疲劳则可以承担相反方向的作用：

\[
J_C
\]

长期过高时抑制进一步认知活动。

因此：

\begin{statementbox}
\centering
\adjustbox{max width=.96\linewidth}{\(\displaystyle
Boredom
\rightarrow
Activity\uparrow,
\)}
\end{statementbox}

而：

\begin{statementbox}
\centering
\adjustbox{max width=.96\linewidth}{\(\displaystyle
Fatigue
\rightarrow
Activity\downarrow.
\)}
\end{statementbox}

两者形成一个双侧稳态机制。

可以写成：

\[
J_C<J_{\min}
\Rightarrow
B\uparrow,
\]

\[
J_C>J_{\max}
\Rightarrow
F\uparrow,
\]

其中 $F$ 表示疲劳状态。

于是：

\[
\boxed{
J_{\min}
<
J_C
<
J_{\max}
}
\]

成为系统比较稳定的认知活动区间。

这说明 AHC 不仅管理：

能量；

威胁；

奖励；

压力，

还可能管理一个更加一般的变量：

\begin{statementbox}
\centering
\adjustbox{max width=.96\linewidth}{\(\displaystyle
\textbf{
认知结构活动强度。
}
\)}
\end{statementbox}

\section{无聊不能直接等于高风险探索}

必须特别强调：

\begin{statementbox}
\centering
\adjustbox{max width=.96\linewidth}{\(\displaystyle
Boredom
\neq
UnboundedExploration.
\)}
\end{statementbox}

如果系统简单采用：

\[
B\uparrow
\Rightarrow
ActionRandomness\uparrow,
\]

则可能产生危险行为。

本章对无聊的定义仅意味着：

\begin{statementbox}
\centering
\adjustbox{max width=.96\linewidth}{\(\displaystyle
\text{认知结构活动需求上升。}
\)}
\end{statementbox}

它并不天然具有：

破坏约束；

越过风险边界；

无限增加行动强度

的权力。

因此：

\[
B
\]

应该改变认知活动倾向，而不是直接指定具体现实行为。

这一区分对于持续智能十分重要：

\begin{statementbox}
\centering
\adjustbox{max width=.96\linewidth}{\(\displaystyle
Drive
\neq
ActionAuthority.
\)}
\end{statementbox}

\section{自主遗忘--无聊耦合假说}

至此，可以把整个推导凝结成一个新的理论假说：

\begin{statementbox}
\centering
\adjustbox{max width=.96\linewidth}{\(\displaystyle
\textbf{
自主遗忘--无聊耦合假说
}
\)}
\end{statementbox}

其内容为：

> 对具有自主遗忘、情景蒸发和长期结构可塑性的持续智能系统，如果长期低认知活动会导致记忆结构净损失，则系统需要某种检测认知结构流量不足并主动恢复认知活动的内稳态机制。无聊可以作为这一机制的一种功能实现。

其因果链可以写成：

\begin{statementbox}
\centering
\adjustbox{max width=.96\linewidth}{\(\displaystyle
AutonomousForgetting
\)}
\end{statementbox}

\[
\Downarrow
\]

\begin{statementbox}
\centering
\adjustbox{max width=.96\linewidth}{\(\displaystyle
PotentialMemoryDrying
\)}
\end{statementbox}

\[
\Downarrow
\]

\begin{statementbox}
\centering
\adjustbox{max width=.96\linewidth}{\(\displaystyle
LowCognitiveStructureFlux
\)}
\end{statementbox}

\[
\Downarrow
\]

\begin{statementbox}
\centering
\adjustbox{max width=.96\linewidth}{\(\displaystyle
BoredomDrive
\)}
\end{statementbox}

\[
\Downarrow
\]

\begin{statementbox}
\centering
\adjustbox{max width=.96\linewidth}{\(\displaystyle
Recall/Recombination/Exploration
\)}
\end{statementbox}

\[
\Downarrow
\]

\begin{statementbox}
\centering
\adjustbox{max width=.96\linewidth}{\(\displaystyle
CognitiveStructureReplenishment.
\)}
\end{statementbox}

这一假说并不意味着任何具有遗忘机制的系统都必然具有类似人类无聊的体验。

更加严格的结论是：

\begin{statementbox}
\centering
\adjustbox{max width=.96\linewidth}{\(\displaystyle
\text{如果自主遗忘造成长期结构损失风险，
则某种功能等价的认知补偿机制具有明显工程必要性。}
\)}
\end{statementbox}

\section{三相记忆与 AHC 的新闭环}

把本章机制重新连接到三相记忆，可以得到：

\[
h
\rightarrow
E
\rightarrow
\Theta.
\]

情景经历形成结构：

\[
E\rightarrow\Theta.
\]

结构成熟以后：

\[
\Theta\rightarrow h',
\]

使预测误差下降：

\[
\varepsilon\downarrow.
\]

于是：

\[
Recall(E)\downarrow,
\]

进一步：

\[
EpisodeReinforcement\downarrow.
\]

最终：

\[
E_{\mathrm{routine}}
\rightarrow
Evaporation.
\]

若新经历长期不足：

\[
J_C\downarrow.
\]

AHC 中无聊状态开始增加：

\[
B\uparrow.
\]

继而推动：

\[
Recall,
\quad
Recombination,
\quad
Exploration
\uparrow.
\]

产生新的当前状态：

\[
h'.
\]

于是形成：

\[
\boxed{
h
\rightarrow
E
\rightarrow
\Theta
\rightarrow
h'
\rightarrow
B
\rightarrow
h''.
}
\]

这里 $B$ 并不是第四种记忆。

它是：

\begin{statementbox}
\centering
\adjustbox{max width=.96\linewidth}{\(\displaystyle
\text{维持三相记忆长期动力循环的内稳态调节变量。}
\)}
\end{statementbox}

\section{理论预测与可检验性}

本章理论产生若干可以进一步实验检验的预测。

第一，如果一个具有情景蒸发机制的持续学习系统长期处于低变化环境，同时禁止内部召回和自主状态生成，则其情景记忆总量应出现持续下降：

\[
\frac{d|E|}{dt}<0.
\]

第二，如果慢结构也具有使用依赖衰减，则长期低活动还应导致部分结构能力下降：

\[
Performance_{\mathrm{rare}}
\downarrow.
\]

第三，如果增加一个基于低认知结构流量的无聊调节机制，则系统应提高：

\[
RecallRate,
\]

\[
InternalActivity,
\]

或：

\[
ExplorationRate.
\]

第四，与没有此机制的系统相比，具有适当无聊反馈的系统在重新进入丰富环境时，应表现出：

更好的旧能力恢复；

更高的状态多样性；

更低的长期结构损失。

第五，如果无聊增益过高，系统可能产生过度探索；如果过低，则无法有效缓解认知结构风干。

因此理论预测存在一个合理控制区间：

\[
\boxed{
B_{\min}
<
B_{\mathrm{gain}}
<
B_{\max}.
}
\]

这些预测使 ``无聊'' 从一个纯粹的拟人化概念转化为可以进行工程实验和动力学分析的功能变量。

\section{理论边界}

本章需要明确区分已有事实、结构类比和理论假说。

第一，有限记忆系统需要遗忘、压缩和资源再利用，是一般工程事实。

第二，长期缺乏使用可能导致某些可塑系统中的功能弱化，以及活动与学习之间存在使用依赖关系，可以作为本理论的重要外部启发。

第三，本章提出的：

\begin{statementbox}
\centering
\adjustbox{max width=.96\linewidth}{\(\displaystyle
\text{自主遗忘}
\rightarrow
\text{认知结构风干风险}
\rightarrow
\text{无聊稳态信号}
\)}
\end{statementbox}

是 AgentOS 理论内部进一步提出的功能假说。

它不等价于对人类主观无聊体验神经机制的完整解释，也不意味着所有智能系统都必须采用同一种 ``无聊'' 变量。

更加一般的理论要求是：

\begin{statementbox}
\centering
\adjustbox{max width=.96\linewidth}{\(\displaystyle
\textbf{
一个具有持续遗忘和长期可塑性的智能系统，需要某种机制避免认知结构活动长期低于维持自身所需的水平。
}
\)}
\end{statementbox}

无聊只是这一功能的一种自然实现。

\section{本章结论：因为智能会遗忘，所以它必须持续产生经历}

三相记忆不是永久存储系统。

工作状态不断消散：

\[
h\rightarrow\varnothing.
\]

情景经历不断蒸发：

\[
E_i\rightarrow\varnothing.
\]

长期结构也可能缓慢变化：

\[
\Theta_t\rightarrow\Theta_{t+1}.
\]

正因为记忆不是静态的，持续智能必须不断产生新的结构活动，以维持长期认知能力。

因此，认知系统中存在两个方向相反的基本流：

\begin{statementbox}
\centering
\adjustbox{max width=.96\linewidth}{\(\displaystyle
ForgettingFlow
\)}
\end{statementbox}

与：

\begin{statementbox}
\centering
\adjustbox{max width=.96\linewidth}{\(\displaystyle
Experience/ReactivationFlow.
\)}
\end{statementbox}

长期稳定要求：

\[
\boxed{
J_{\mathrm{experience}}
+
J_{\mathrm{reactivation}}
\approx
J_{\mathrm{forgetting}}.
}
\]

当左侧长期不足时，系统出现认知结构风干风险。

无聊可以被理解为检测这种失衡并推动系统恢复认知活动的一种 AHC 稳态变量。

因此，本章最终提出：

\begin{statementbox}
\centering
\adjustbox{max width=.96\linewidth}{\(\displaystyle
\textbf{
无聊并不首先来自“没有事情可做”，
而来自一个具有自主遗忘能力的持续智能系统对长期认知结构流量不足的检测。
}
\)}
\end{statementbox}

进一步：

\begin{statementbox}
\centering
\adjustbox{max width=.96\linewidth}{\(\displaystyle
\textbf{
自主遗忘使持续智能能够避免历史无限堆积；
无聊则可能使持续智能避免因为遗忘过度而失去自身。
}
\)}
\end{statementbox}

两者共同形成：

\begin{statementbox}
\centering
\adjustbox{max width=.96\linewidth}{\(\displaystyle
\textbf{
遗忘负责释放过去，
无聊负责重新寻找未来。
}
\)}
\end{statementbox}

而从三相记忆动力学的角度，可以将本章最核心的结论进一步压缩为：

\begin{statementbox}
\centering
\adjustbox{max width=.96\linewidth}{\(\displaystyle
\textbf{
因为智能会遗忘，所以它必须拥有继续经历世界的内生理由。
}
\)}
\end{statementbox}